\documentclass[aspectratio=43]{beamer}

\usepackage[T1]{fontenc}
\usepackage[utf8]{inputenc}
\usepackage[english]{babel}
\usepackage{lmodern}
\usepackage{booktabs}
\usepackage{multirow}
\usepackage{subcaption}

%import the theme and packages for the code blocks
\input{code_blocks_theme.tex}

% Use Unipd as theme, with options:
% - pageofpages: define the separation symbol of the footer page of pages (e.g.: of, di, /, default: of)
% - logo: position another logo near the Unipd logo in the title page (e.g. department logo), passing the second logo path as option 
% Use the environment lastframe to add the endframe text
\usetheme[pageofpages=of]{Unipd}

\title{A Comparative Study of Image Inpainting Methods
	under Diverse Masking Conditions\\}
\subtitle{}
\author[]{Ponchio Mattia, Secco Benedetto}

\date{15th September 2026}


\begin{document}
	\frame{\titlepage}
	\section{Partial Convolution}
	\begin{frame}
		\begin{center}
			\Huge{\textbf{Partial Convolution}}
		\end{center}
	\end{frame}
	
	
	\begin{frame}{Goals}
		\begin{itemize}
			\item Independent from hole initialization values.
			\item No expensive post-processing needed.
			\item Convolution results depend only on the non-hole regions at every layer, with the masked region being automatically updated at each step.
			\item Handle both regular and irregular shaped holes, regardless of their proximity to the borders of the image.
		\end{itemize}
	\end{frame}
	
	
	\begin{frame}{Partial Convolutional Layer}
		\begin{center}			
			\begin{equation*}
				x' = \begin{cases}
					\mathbf{W}^T(\mathbf{X}\odot\mathbf{M})\frac{\text{sum}(\mathbf{1})}{\text{sum}(\mathbf{M})} + b,  & \text{if }  \text{sum}(\mathbf{M})>0  \\
					0,  & \text{otherwise}
				\end{cases}
			\end{equation*}
			\begin{itemize}
				\item Mask convention: 0 for masked portion.
				\item Output values depend only on unmasked inputs.
				\item Scaling factor to adjust for the varying amount of unmasked inputs.
				\item If the convolution was able to condition its output on at least one valid input value, then that location becomes valid, expressed with:
			\end{itemize}
			\begin{equation*}
					m' = \begin{cases}
					1,  & \text{if }  \text{sum}(\mathbf{M})>0  \\
					0,  & \text{otherwise}
				\end{cases}
			\end{equation*}
		\end{center}
	\end{frame}
	
	
	\begin{frame}{Network Architecture Inspiration}
		\centering
		\includegraphics[scale=0.14]{images/u-net-architecture.png}
		\begin{itemize}
			\item UNet-like architecture: Encoder-Decoder with skip connections.
		\end{itemize}
		\small
		\footnote{Image source: Ronneberger, O., Fischer, P., Brox, T.: U-net: Convolutional networks for biomedical image segmentation. Springer (2015)}
	\end{frame}
	
	
	\begin{frame}{Network Architecture Details}
		\begin{itemize}
			\item Each convolution is transformed into a partial convolution.
			\item The skip links concatenate two feature maps and two masks respectively, acting as the feature and mask inputs for the next partial convolution layer.
			\item The last partial convolution layer’s input contains the concatenation of the original input image with hole and original mask, making it possible for the model to copy non-hole pixels.
		\end{itemize}
	\end{frame}
	
	
	\begin{frame}{Network Architecture Visualization}
		\centering
		\includegraphics[width=0.9\textwidth, height=0.85\textheight]{images/Pconv_net.png}
	\end{frame}
	
	
	\begin{frame}{Loss Function}
		\centering
		\begin{equation*}
				\mathcal{L}_{total} = 
				\mathcal{L}_{valid} + 6\mathcal{L}_{hole} + 0.05\mathcal{L}_{perceptual}
				+ 120(\mathcal{L}_{style_{out}} + \mathcal{L}_{style_{comp}}) + 0.1\mathcal{L}_{tv}
		\end{equation*}
		\begin{itemize}
			\item Multiple terms for taking into account both pixel-wise accuracy and perceptual metrics.
			\item $\mathcal{L}_{hole}$ and $\mathcal{L}_{valid}$ are the $L_1$ losses on the network output for the hole and the non-hole pixels respectively.
		\end{itemize}
		\begin{equation*}
			\mathcal{L}_{perceptual} = \sum_{p=0}^{P-1} \frac{\| \Psi_{p}^{\mathbf{I}_{out}} - \Psi_{p}^{\mathbf{I}_{gt}} \|_1}{N_{\Psi_{p}^{\mathbf{I}_{gt}}}}
			+ \sum_{p=0}^{P-1} \frac{\|  \Psi_{p}^{\mathbf{I}_{comp}} - \Psi_{p}^{\mathbf{I}_{gt}} \|_1}{N_{\Psi_{p}^{\mathbf{I}_{gt}}}}
		\end{equation*}
		\begin{itemize}
			\item $\mathbf{I}_{comp}$ is the raw output image $\mathbf{I}_{out}$, but with the non-hole pixels directly set to ground truth.
			\item The perceptual loss computes the $L_1$ distances between both $\mathbf{I}_{out}$ and $\mathbf{I}_{comp}$ and the ground truth, but after projecting these images into higher level feature spaces using an ImageNet-pretrained VGG-16.
		\end{itemize}
	\end{frame}
	
	
	\begin{frame}{Loss function}
		\centering
		\begin{itemize}
			\item The style-loss term is similar to the perceptual loss, but an autocorrelation (Gram matrix) is performed on each feature map before applying the $L_1$.
			\item As before, $\mathcal{L}_{style_{comp}}$ uses the raw output image but with non-hole pixels set to ground truth.
			\item The total variation loss $\mathcal{L}_{tv}$ represents the smoothing penalty on $R$, where $R$ is the region of 1-pixel dilation of the hole region, still evaluated using the $L_1$ norm.
		\end{itemize}
	\end{frame}
	
	
	\begin{frame}{Results}
		content...
	\end{frame}
	
	
	\section{Contextual Attention}
	\begin{frame}
		\begin{center}
			\Huge{\textbf{Contextual Attention}}
		\end{center}
	\end{frame}
	
	
	\begin{frame}{Goals}
		\centering
		\begin{itemize}
			\item Combine the ability of deep learning based approaches to generate visually plausible inpainting with patch synthesis methods's to borrow textures from the surrounding regions.
			\item Deal with multiple holes at arbitrary locations and with variable sizes.
			\item Integrate attention mechanism in the network: learn where to borrow or copy feature information from known background patches to generate missing patches.
		\end{itemize}
	\end{frame}
	
	
	\begin{frame}{Contextual Attention Layer}
		\centering
		\includegraphics[width=\textwidth, height=0.5\textheight]{images/contextual_attention_layer.jpg}
		\begin{itemize}
			\item Match features of missing pixels (foreground) to surroundings (background) via cosine similarity $\rightarrow$ weigh it with scaled softmax to get attention score for each pixel.
			\item Further encourage coherency of attention by propagation (fusion) $\rightarrow$ left-right propagation followed by a top-down propagation with kernel size of k.
		\end{itemize}
	\end{frame}
	
	
	\begin{frame}{Network Architecture: $1^{st}$ stage}
		\centering
		\includegraphics[width=\textwidth, height=0.35\textheight]{images/framework.jpg}
		\begin{itemize}
			\item Two-stage coarse-to-fine network architecture where the first network makes an initial coarse prediction, and the second network takes the coarse prediction as inputs and predict refined results.
			\item The coarse network is a fully convolutional encoder-decoder architecture with dilated convolutions.
			\item It is trained with the reconstruction loss explicitly; no $BN$ layers and ELU as activation.
		\end{itemize}
	\end{frame}
	
	
	\begin{frame}{Spatially Discounted Loss}
		\centering
		\begin{itemize}
			\item Since the original image is the only ground truth to compute a reconstruction loss, strong enforcement of reconstruction loss in masked pixels may mislead the training process of the convolutional network.
			\item Missing pixels near the hole boundaries have much less ambiguity than those pixels closer to the center of the hole.
			\begin{center}
				$\downarrow$
			\end{center}
			\item Weight mask \(\mathbf{M}\) where the weight of each pixel in the mask is computed as \(\gamma^{\ell}\), where \(\ell\) is the distance of the pixel to the nearest known pixel.
		\end{itemize}
	\end{frame}
	
	
	\begin{frame}{Network Architecture: $2^{nd}$ stage}
		\centering
		\includegraphics[width=\textwidth, height=0.35\textheight]{images/unified_attention_network.jpg}
		\begin{itemize}
			\item Two parallel encoders: the bottom encoder specifically focuses on hallucinating contents with layer-by-layer (dilated) convolution, while the top one tries to attend on background features of interest.
			\item Output features from two encoders are aggregated and fed into a single decoder to obtain the final output.
			\item Color to indicate the relative location of the most interested background patch for each foreground pixel: white means the pixel attends on itself, pink on bottom-left, green on top-right.
		\end{itemize}
	\end{frame}
	
	
	\begin{frame}{$2^{nd}$ stage Loss}
		\centering
		\begin{itemize}
			\item Reconstruction loss and WGAN-GP loss (on both global and local outputs
			of the second-stage refinement network to enforce global and local consistency).
			\item The pixel-wise reconstruction loss directly regresses
			holes to the current ground truth image, while WGANs implicitly learn to match potentially correct images and train the generator with adversarial gradients.
			\item The gradient penalty (applied only to pixels inside the holes) contributes to preventing mode collapse and balances the two critics.
			\item No perceptual, style or total variation losses used, as they do not bring noticeable improvements in this case.
			\begin{equation*}
				\begin{split}
					\min_G \max_{D \in \mathcal{D}} & \mathbb{E}_{\mathbf{x} \sim \mathbb{P}_r}[D(\mathbf{x})] - \mathbb{E}_{\tilde{\mathbf{x}} \sim \mathbb{P}_g}[D(\tilde{\mathbf{x}})] \\
					& + \lambda \mathbb{E}_{\hat{\mathbf{x}} \sim \mathbb{P}_{\hat{\mathbf{x}}}} \left[ \left( \lVert \nabla_{\hat{\mathbf{x}}} D(\hat{\mathbf{x}}) \odot (\mathbf{1} - \mathbf{m}) \rVert_2 - 1 \right)^2 \right]
				\end{split}
			\end{equation*}
		\end{itemize}
	\end{frame}
	
	
	\begin{frame}{Results}
		content...
	\end{frame}
	
	
\end{document}